\section{总结与展望}
\label{sec:summary}
我们发展了在小流时极限下从梯度流方案匹配到 $\overline{\rm MS}$ 方案、适用于威尔逊线算符的系统方法。该方法可推广到多种应用，包括准PDF、pNRQCD 框架下夸克偶素衰变与产生中出现的胶子关联函数，以及以色磁场插入威尔逊圈表述的自旋依赖势。

在一维辅助场形式中，本工作所讨论的威尔逊线算符从梯度流方案到 $\overline{\rm MS}$ 方案的匹配，在小流时极限下可以约化为对相应局域流算符的匹配。
我们通过比较这两种方案下的微扰计算，把这些局域流算符从梯度流方案匹配到 $\overline{\rm MS}$ 方案。
利用小流时展开并尽可能把外动量取为零，可以使匹配过程得到极大的简化。
这样便可以应用 HQET 中成熟的匹配计算方法。完成局域流算符的匹配之后，小流时极限下相应威尔逊线算符的匹配系数便自动得到，
因为由于威尔逊线算符的可乘重正化性，不同时空点上重正化局域流算符的组合除减除掉的线性发散外不再引入更多的 UV 发散。


由于局域流算符的可乘重正化性，局域流算符的匹配还可以进一步分解为场与顶点的匹配。
于是我们计算了辅助``重''辅助场 $h_v$、无质量夸克场 $\psi$、$\psi h_v$ 顶点、胶子场 $gA$ 以及局域流算符 $gF^{\mu\nu}_{\|\perp}h_v,\, gF^{\mu\nu}_{\perp\perp}h_v$ 与 $g\bar{h}_vF^{\mu\nu}_{\perp\perp}h_v$ 的顶点在单圈水平的匹配系数。我们还得到了线性发散的``质量''修正 $\delta m$ 的单圈结果。有了这些结果，我们给出了小流时极限下局域流算符与威尔逊线算符的匹配系数。

值得一提的是，该方法也适用于在小格距极限下从格点正则化方案到 $\overline{\rm MS}$ 方案的匹配，只是由于 UV 与 IR 发散的处理而带有额外的复杂性 \cite{Constantinou:2017sej}。



在附录中我们证明：对于 eq.~(\ref{eq:quarkhmte}) 定义的零外动量强子矩阵元，在单圈水平上，只要流半径小于物理距离 $z$——格点梯度流计算中通常满足这一条件——有限流时效应就非常小。


利用我们为威尔逊线算符发展的从梯度流方案匹配到 $\overline{\rm MS}$ 方案的系统方法，
把当前研究推广到两圈阶是直截了当的，这将更有助于用非光锥威尔逊线算符构造的矩阵元的格点梯度流计算。我们把它留作未来的研究。

\subsubsection*{补充说明}
在本工作提交之前，我们注意到胶子关联函数 $G_B$ 匹配系数的单圈结果最近发表在文献 \cite{Altenkort:2023eav} 中，它与我们的 $c_F^2(t, \mu)$ 单圈结果相同，因为胶子关联函数 $G_B$ 的匹配计算可以约化为对局域流算符 $g\bar{h}_vF^{\mu\nu}_{\perp\perp}h_v$ 的匹配计算。
